\chapter{Background and Related Work}
\label{ch:background}

\section{Spherical harmonics as a regression basis}
\label{sec:sh-background}
 
Let $S^2 = \{(\theta,\phi) \mid \theta \in [0,\pi],\ \phi \in [0,2\pi)\}$ denote the unit sphere where colatitude ($\theta$) and longitude ($\phi$). The (complex) spherical harmonics of degree $l \in \mathbb{N}_0$ and order $m \in \{-l,\dots,l\}$ are defined as
\begin{equation}
    Y_l^m(\theta,\phi) = \sqrt{\frac{2l+1}{4\pi}\frac{(l-m)!}{(l+m)!}}\, P_l^m(\cos\theta)\, e^{im\phi},
    \label{eq:sph-harm}
\end{equation}
where $P_l^m$ is the associated Legendre function \citep{thomas2025}. These functions arise from separating variables in the Helmholtz eigenvalue problem $\nabla^2\Psi = \lambda\Psi$ on $\mathbb{R}^3$, restricted to a sphere of fixed radius; writing $\Psi(r,\theta,\phi)=R(r)\Theta(\theta)\Phi(\phi)$ and separating variables yields the associated Legendre equation for $\Theta$ and a trivial harmonic equation for $\Phi$, whose periodic solutions force the order $m$ to be an integer. (Given the full derivation in the predecessor dissertation's Appendix A.)
 
Two properties are particularly important for the spherical harmonics : orthogonality and Laplace-Beltrami eigenfunction property. First, orthonormality on the sphere with respect to the $L^2$ inner product,
\begin{equation}
    \langle Y_l^m, Y_{l'}^{m'}\rangle = \int_0^{2\pi}\!\!\int_0^\pi Y_l^m(\theta,\phi)\,\overline{Y_{l'}^{m'}(\theta,\phi)}\,\sin\theta\,\mathrm{d}\theta\,\mathrm{d}\phi = \delta_{ll'}\delta_{mm'},
    \label{eq:orthogonality}
\end{equation}
means that any function in $L^2(S^2)$ has a unique expansion in this basis, and that the squared $L^2$ norm of a truncated expansion is simply the sum of squared coefficient magnitudes, $\|Sf\|^2_{L^2(S^2)} = \sum_{l,m}|v_{lm}|^2$ , which allows the regularisation term in \cref{eq:reg-lsq} below to be written purely in terms of the coefficient vector $v$ rather than as an explicit integral, so that regularisation can be expressed directly in terms of the harmonic coefficients, without having to numerically evaluate an integral for the penalty. Second, spherical harmonics are eigenfunctions of the spherical Laplace-Beltrami operator: $\nabla^2_\Omega Y_l^m = -l(l+1)Y_l^m$, where the eigenvalue depends only on the degree $l$. Low $l$ leads to smooth patterns across the sphere, whereas high $l$ leads to fine detailed patterns. Thus, high-degree harmonics can represent fine scale changes, but they might also cause noise in the data. Higher $l$ values correspond to larger eigenvalues of the spherical harmonic operator, allowing high-degree components to be penalised more strongly. This allows us to control the smoothness of the reconstructed field directly through the spherical harmonic coefficients. 

The smoothness penalty is measured using \textbf{Sobolev spaces}. In general, Sobolev spaces contain functions whose weak derivatives up to order $k$ are sufficiently well behaved in $L^p(S^2)$. The Hilbert-space case, with p=2, is written as $H^k(S^2):=W^{k,2}(S^2)$ and the corresponding Sobolev norm $\|f\|_{H^k}^2 = \sum_{|\alpha|\le k}\|D^\alpha f\|^2_{L^2}$, is
with relation to Sobolev seminorm $|f|_{H^k}^2 = \|D^k f\|^2_{L^2}$ . This seminorm focuses only on the highest-order derivative, rather than all derivatives up to order $k$. Using the spherical-harmonic eigenvalue relationship, the Sobolev seminorm of the truncated expansion in \cref{eq:truncated-expansion} can be written as $|Sf|^2_{H^k(S^2)} = \sum_{l,m}(l(l+1))^k|v_{lm}|^2$. This expression shows how the smoothness penalty acts on different spherical-harmonic degrees, where low-degree coefficients are penalised less and high-degree coefficients are penalised more. Therefore, a field with mostly low-degree energy is smoother and will have a smaller Sobolev seminorm. In contrast, a field with substantial high-degree energy, containing more detailed high-frequency structure, has a larger Sobolev seminorm. Penalising $|Sf|^2_{H^k(S^2)}$ in the least-squares objective provides a direct and interpretable way of favouring smoother reconstructions. The penalty here discourages excessive high-frequency components, and the order $k$ controls how strongly the high-degree/high-frequency structure is penalised. Therefore, larger $k$ means more aggressive suppression of high-frequency components, while smaller $k$ corresponds to less aggressive smoothness control. 
 
Given $n$ noisy observations $\{\hat f_i\}_{i=1}^n$ of an unknown scalar field $f(\theta,\phi)$ collected at spherical locations $\{(\theta_i,\phi_i)\}_{i=1}^n \subset S^2$, the unknown field is approximated by using a truncated spherical-harmonic expansion:
\begin{equation}
    Sf(\theta,\phi) = \sum_{l=0}^{L}\sum_{m=-l}^{l} v_{lm}\, Y_l^m(\theta,\phi),
    \label{eq:truncated-expansion}
\end{equation}
where $L$ controls model complexity and $v_{lm} \in \mathbb{C}$ are unknown harmonic coefficients. The design matrix is written as $A \in \mathbb{C}^{n \times (L+1)^2}$ with entries $A_{i,(l,m)} = Y_l^m(\theta_i,\phi_i)$,  contains the spherical-harmonic basis functions evaluated at the observation locations. The coefficient vector $v$ contains all $v_{lm}$ values. The number of coefficients is $(l+1)^2$ , and the coefficients are estimated by solving a regularised least-squares problem
\begin{equation}
    \min_{v \in \mathbb{C}^{(L+1)^2}} \; \|Av - \hat f\|_2^2 + \lambda\, |Sf|^2_{H^k(S^2)},
    \label{eq:reg-lsq}
\end{equation}
where the first term measures the difference between the reconstructed values and the observed data, and the second term controls the smoothness of the reconstruction. $\lambda$ controls the strength of the regularisation and $k$ determines how the Sobolev penalty weights different spherical-harmonic degrees. The Sobolev seminorm can be written directly using the spherical-harmonic coefficients $|Sf|^2_{H^k(S^2)} = \sum_{l,m} (l(l+1))^k |v_{lm}|^2$. This is possible because of the orthogonality/eigenvalue properties of the spherical-harmonic basis. The penalty therefore acts directly on the coefficients. 

Differentiating the regularised objective with respect to $v$ gives the regularised normal equation. Since the spherical-harmonic coefficients are complex, the \textbf{Hermitian inner product}, regularised normal equation is used:
\begin{equation}
    v = \left(A^H A + \lambda M_k\right)^{-1} A^H \hat f,
    \label{eq:normal-eq}
\end{equation}
where $A^H$ denotes the Hermitian transpose of $A$ and $M_k$ is the diagonal matrix with entries $(l(l+1))^k$ on the diagonal, with one for each $(l,m)$ pair. Setting $k=0$, the method becomes ordinary Tikhonov/ridge regularisation. Increasing $k$ gives greater penalties to high-degree spherical-harmonic components, therefore penalises high-frequency, rough components of the reconstruction. This is suitable when the underlying field is expected to be relatively smooth. For example, large-scale surface temperature is expected to contain mainly relatively low-order spatial structure, while local observations can contain higher-frequency/noisy variation. The regularisation therefore balances between fitting the observed data and maintaining a smooth reconstructed field. 

The family of regularisers can be viewed as a spectral extension of ordinary Tikhonov regularisation. In this regression, the single scalar penalty weight $\lambda$ of classical ridge regression is replaced by a degree-dependent weight $\lambda\,(l(l+1))^k$. The penalty increases polynomially with $l$ according to parameter $k$. From a bias-variance perspective, increasing $\lambda$ at fixed $k$ produces stronger regularisation overally, and reduced variances. This cause the fitted coefficients to be less sensitive to the particular noisy sample observed. However, it increases bias, as the fitted model is pulled towards zero except for the low-degree, unpenalised structure. This means that $\lambda$ controls the overall strength of the bias-variance trade-off. In contrast, increasing $k$ while keeping $\lambda$ fixed does not apply the same additional penalty to all coefficients. Instead, it increases the shrinkage more strongly for high-degree coefficients, while low-degree coefficients representing large-scale structure are affected less. Therefore, $\lambda$ controls the overall strength of regularisation, whereas $k$ determines how the regularisation is distributed across different degrees. This differs from ordinary ridge regression, where the same penalty is applied to all coefficients without distinguishing between low-degree, large-scale components and high-degree, small-scale components. The Sobolev-based regulariser instead provides more targeted shrinkage, allowing high-degree components to be suppressed more strongly while preserving important low-degree structure. This is particularly useful for smooth spatial fields. This dissertation's own independent grid search consistently favours $k>0$ over $k=0$ for climate-type fields which follows the finding in predecessor dissertation. This is consistent with the expected characteristics of climate-type fields, whose genuine physical structure is concentrated at low-to-moderate degree. At high $l$, variation may contain observational noise, small-scale spatial effects or station-specific effects,  rather than large-scale climate signal. Therefore, stronger penalisation of high-degree coefficients is physically reasonable. This provides an explanation for why the model-selection results favour $k>0$ rather than ordinary ridge regularisation with $k=0$.

This is exactly the model developed and validated on synthetic and real climate/wind data in the predecessor dissertation \citep{thomas2025}. First, the predecessor dissertation found an approximate relationship between the number of observations $n$, and the optimal spherical-harmonic truncation degree $L$. For uniformly distributed observations $L \propto n^{1/2}$. This relationship was observed for achieving a fixed target reconstruction accuracy. Therefore, as the number of observations increases, a higher model degree can be used while maintaining the desired reconstruction accuracy. Second, the previous work also investigated regularisation using the second Sobolev seminorm $k=2$. The choice of $k=2$ was motivated by the expected connection between the physical field and a diffusion/heat-type equation. The results showed that regularisation with $k=2$ improved reconstruction accuracy compared with unregularised least squares. The improvement was particularly noticeable when observations were sparse and the desired model complexity was relatively high. This dissertation adopts \cref{eq:reg-lsq}--\cref{eq:normal-eq} unchanged as its base predictor, the contribution here lies entirely in what is layered on top of it.
 
\section{Split conformal prediction}
\label{sec:cp-background}

Conformal prediction \citep{vovk2005} is a general framework for constructing prediction regions with a finite-sample, distribution-free coverage guarantee. It requires only that the calibration and test data be \emph{exchangeable} defined as their joint distribution is invariant under permutation of the indices, and no assumption on the underlying data distribution or on the correctness of the base predictor is required. This is a substantially easier requirement to achieve than the independent-and-identically-distributed assumption typically invoked in statistical learning theory. i.i.d.\ data is exchangeable, but exchangeable data do not need to be i.i.d.; thus, the framework has proved attractive well beyond its original setting. The originally proposed full transductive version of conformal prediction requires refitting the base predictor once for every label at every test point. This is computationally prohibitive for anything beyond small models and discrete label spaces.

In its simplest \emph{split} (or inductive) form \citep{lei2018} sacrifices some statistical efficiency for a large reduction in computational cost and is the form used throughout this dissertation. The available data is partitioned once into a training set $\mathcal{D}_{tr}$ to fit an arbitrary base predictor $\hat y$, and a calibration set $\mathcal{D}_{cal}$ to calibrate the width of the prediction region. This process uses the fixed base predictor without refitting it afterward. For each calibration point, a \emph{nonconformity score} $s(X_i,Y_i)$ is computed in the simplest regression case which is the absolute residual $|Y_i - \hat y(X_i)|$.  The $\lceil (1-\alpha)(n_{cal}+1)\rceil / n_{cal}$ empirical quantile of these scores, $\hat q_\alpha$ is taken as the calibrated threshold, then the prediction region at a new point $x$ will be $C(x) = \{y : s(x,y) \le \hat q_\alpha\}$. The symmetric interval is constructed as $[\hat y(x) - \hat q_\alpha,\ \hat y(x) + \hat q_\alpha]$ for the absolute-residual score. By the exchangeability of $\mathcal{D}_{cal}$ and a held-out test point, it guarantees
\begin{equation}
    \mathbb{P}\big(Y_{n+1} \in C(X_{n+1})\big) \ge 1 - \alpha
\end{equation}
marginally over the joint distribution of $(X,Y)$, regardless of the choice of nonconformity score or the performance of fitting the base predictor. The reasoning behind this guarantee is straightforward and helps clarify the assumptions required. Conditional on the training set $\mathcal{D}_{tr}$, the fitted prediction function $\hat y$ is fixed. Since $\mathcal{D}_{cal}\cup\{(X_{n+1},Y_{n+1})\}$ is exchangeable by assumption, the corresponding nonconformity scores ${s(X_i,Y_i)}{i\in\mathcal{D}{cal}}$ with the test score $s(X_{n+1},Y_{n+1})$ are also treated as exchangeable. This forms a sequence of $n_{cal}+1$ real numbers because $s(\cdot,\cdot)$ is a fixed function once $\hat y$ has been fitted, and it is then applied to each calibration and test pair. So the rank of the last element is uniformly distributed over ${1,\dots,n_{cal}+1}$for any exchangeable sequence of real numbers. Thus, the probability that test score is larger than the $\lceil(1-\alpha)(n_{cal}+1)\rceil$-th order statistic of calibration scores is at most $\alpha$, which is equivalent to $Y_{n+1}\notin C(X_{n+1})$. This statement does not depend on the specific geometry used to measure ``closeness'', or on how well $\hat y$ predicts. A weak predictor will still provide valid marginal coverage, even though its prediction intervals may be wider. The genuinely distribution-free property allows the score function to vary for geometric reasons, as implemented in \cref{sec:geodesic-cp-background,sec:adaptation}, without needing to re-derive the coverage argument from the beginning.

However, marginal coverage is only an overall guarantee. It measures coverage across the entire test set, but it does not guarantee equal or reliable coverage across different regions of the input space. For example, some regions may have higher coverage while others have lower coverage. This is especially important when prediction difficulty varies across locations. This matters for climate reconstruction, because climate reconstruction data can have heterogeneous prediction errors. The residual distribution may be heteroscedastic, meaning the error variability changes with location. For example, densely sampled temperate regions are generally easier to predict, while sparsely sampled regions are potentially harder to predict. Additionally, physically unusual regions, such as high-altitude sites could potentially lead to higher prediction errors. In this case, the problem comes when using standard conformal prediction because standard marginal coverage treats the test set as a whole. It does not explicitly account for location-dependent prediction difficulty. Instead, locally adaptive conformal prediction can account for the local prediction difficulty as a function of the input covariates, by incorporating this estimate into the nonconformity score. Thus, easier locations can receive smaller prediction intervals and harder locations can receive larger prediction intervals. The main motivation is to move beyond simply achieving overall or marginal coverage. This dissertation aims for more consistent coverage across regions with different levels of prediction difficulty. A commonly used formulation is
\begin{equation}
s(x,y) = \frac{|y-\hat y(x)|}{\hat\sigma(x)},
\label{eq:adaptive-score}
\end{equation}
\begin{equation}
C(x) =
\left[
\hat y(x)-\hat q_\alpha\hat\sigma(x),\,
\hat y(x)+\hat q_\alpha\hat\sigma(x)
\right].
\label{eq:adaptive-interval}
\end{equation}

Thus, the interval will wider in regions at higher estimated prediction difficulty. This approaches retain the marginal conformal coverage guarantee if the score is constructed following the required exchangeability conditions. Meanwhile it also potentially improving the allocation of interval width across regions with different levels of difficulty \citep{lei2018,romano2019}.


\section{Geometry-aware and adaptive geodesic conformal prediction}
\label{sec:geodesic-cp-background}

\citet{amirishahbazi2026} extend the locally-adaptive idea to cases where the \emph{response} itself is manifold-valued. For instance, a direction on $S^2$, as in geomagnetic field forecasting, or an angle pair on the flat torus $T^2$, as in protein backbone dihedral angles. They observe that applying ordinary Euclidean conformal prediction is not appropriate for these outputs. One might instead represent the manifold using coordinates, such as longitude and latitude/colatitude, and treat these as ordinary Euclidean coordinates. But Euclidean coordinate distance does not represent the true geometry uniformly across the sphere. This is because a fixed coordinate rectangle with fixed half-width $\delta$ has area $2\delta\,(\cos(\theta-\delta) - \cos(\theta+\delta))$ that produces different intrinsic areas depending on location. It shrinks to zero near the poles as the azimuthal direction is metrically compressed there. Therefore, coordinate-based regions can behave differently depending on location. They therefore use geodesic distance $\dgeo(u,v) = \arccos\langle u,v\rangle$, giving geodesic-cap prediction regions with area $2\pi(1-\cos r)$, which depends on cap radius $r$ and not on position.

They combine the geometric correction with the locally adaptive conformal idea of \citet{lei2018,romano2019} through the \emph{adaptive geodesic} score:
\begin{equation}
    s_{\mathrm{adapt}}(x,y) = \dgeo\big(\hat y(x), y\big) \,/\, \hat\sigma(x),
    \label{eq:paper-score}
\end{equation}
where the difficulty estimator $\hat\sigma(x)$ is produced by a two-stage procedure using only the training set. This preserves the calibration-set exchangeability requirement. Five-fold cross-validated geodesic residuals $e_i = \dgeo(\hat y^{(-k)}(X_i), Y_i)$ are calculated using the training data, while a $k$-nearest-neighbour regressor with $k=20$ is fit to estimate $e_i$ from $X_i$.  The resulting prediction region is then $B_{\mathrm{geo}}(\hat y(x), \hat q_\alpha \hat\sigma(x))$, where easier locations correspond to smaller $\hat{\sigma}(x)$ values, and harder locations require larger $\hat{\sigma}$ values. Their method is then tested on synthetic spherical data, as well as on real geomagnetic-field forecasting derived from the IGRF-14 model \citep{alken2021}. Their method reduces variation in conditional coverage, improves worst-case bin-level coverage, and improves performance relative to both the non-adaptive geodesic method and the naive coordinate-based method.
 
\citet{amirishahbazi2026} situate this contribution within a still relatively sparse literature on conformal prediction for manifold-valued outputs, structured outputs, or multivariate outputs. \citet{kuleshov2018} study the calibration of regression uncertainty estimates and propose a post-hoc recalibration procedure for regression models. Their work is related to the broader objective of obtaining calibrated predictive uncertainty, but it differs in calibration framework and does not specifically address the intrinsic geometry of curved output spaces. \citet{messoudi2021,messoudi2022} address multi-target regression using copula-based and ellipsoidal prediction regions, respectively. However, their approaches remain in Euclidean output spaces and do not deal directly with curved or manifold output geometry. More recent extensions broaden conformal prediction along orthogonal axes to the geometric one addressed here: (1) hierarchical and group-aware calibration has been applied to multiple healthcare data \citep{amirishahbazi2026health}; (2) multi-scale conformal prediction has been developed at different levels of resolution \citep{baheri2025scales}; and (3) split conformal prediction has been extended to function-valued outputs in neural-operator settings \citep{millard2025neural}. These papers show that conformal prediction is being extended from different perspectives of structured data. However, they do not directly address the particular problem of constructing conformal prediction regions using the intrinsic geometry of a curved output space. This is the setting addressed by \citet{amirishahbazi2026}, the most directly relevant previous work, whose approach provides the methodological starting point for the adaptation developed in this dissertation.

\section{The methodological gap}
\label{sec:adaptation}

The three pieces of background above do not compose automatically. \citet{thomas2025}'s spherical harmonic model reconstructs a \emph{scalar} field ($f: S^2 \to \mathbb{R}$). It has an $S^2$-valued \emph{input} but a Euclidean \emph{output}. In contrast, \citet{amirishahbazi2026}'s adaptive geodesic conformal prediction addresses the reverse case, where a possibly Euclidean input is paired with a manifold-valued output, such as a point on $S^2$. The geodesic distance $\dgeo(\hat y(x), y)$ is then defined between the predicted and true outputs on the output manifold. Applying \cref{eq:paper-score} literally to the temperature reconstruction problem is not meaningful, since there is no geodesic distance between two temperature values, $\mathbb{R}$ is flat and already the natural output space.

The adaptation work in this dissertation is to recognise the part of \citet{amirishahbazi2026}'s method that is genuinely about \emph{geometry}. The main idea is to separate two different parts of the original method which are the use of geometry and the fact that the response output is manifold-valued. The geometric idea can still be useful even though temperature is scalar by using geodesic rather than coordinate distance to decide which points are "nearby" for the purposes of estimating local difficulty. Therefore, geodesic distance can be used to determine which observation locations are geographically close to each other. This is then used to estimate local prediction difficulty. Concretely:

\begin{itemize}
    \item \textbf{Nonconformity score}. Because the temperature response is scalar, the ordinary Euclidean locally-adaptive score is used, $s(x,y) = |y-\hat y(x)| / \hat\sigma(x)$, which is exactly the same as \citet{lei2018} and \citet{romano2019}, since the temperature response lives in $\mathbb{R}$. The numerator measures the temperature prediction error, while $\hat\sigma(x)$ represents the estimated difficulty of predicting at location $x$. 

    \item \textbf{Geodesic distance for the input locations}. The sphere's geometry is used when estimating $\hat\sigma(x)$. Instead of calculating nearest neighbours using ordinary Euclidean distance between $(\theta,\phi)$ coordinates, great-circle/geodesic distance is used. This avoids problems caused by treating latitude/longitude as ordinary Euclidean coordinates from exactly the pole-distortion problem \citeauthor{amirishahbazi2026} identified. For each query location, the $k=20$ nearest neighbours are found using the great-circle distance
    \begin{equation}
        \dgeo\big((\theta_1,\phi_1),(\theta_2,\phi_2)\big) = \arccos\big(\sin\theta_1\sin\theta_2\cos(\phi_1-\phi_2) + \cos\theta_1\cos\theta_2\big),
        \label{eq:geo-input}
    \end{equation}and then $\hat\sigma(x)$ is the mean of the absolute, out-of-fold residuals of these $k$ nearest neighbours. Thus, the local difficulty estimate reflects the actual spherical distance between locations. 

    \item  \textbf{Prediction region} at a query location is a varied interval in temperature defined as $[\hat y(x) - \hat q_\alpha \hat\sigma(x),\ \hat y(x) + \hat q_\alpha \hat\sigma(x)]$. It is not a geodesic cap, but its radius is location-dependent via $\hat\sigma(x)$. This means the width changes depending on the estimated local difficulty. This follows the same adaptive principle as the source paper, even though the actual prediction region is different. 
\end{itemize}

\Cref{tab:adaptation-summary} summarises this correspondence precisely, since it is the main methodological contribution in the dissertation.
\begin{table}[H]
\centering
\caption{How this dissertation adapts adaptive geodesic conformal prediction from a manifold-valued-response setting to a scalar-response, spherical-input setting.}
\label{tab:adaptation-summary}
\begin{tabular}{p{3.6cm}p{5.1cm}p{5.1cm}}
\toprule
& \citet{amirishahbazi2026} & This dissertation \\
\midrule
Manifold-valued object & Response $Y \in S^2$ & Covariate $X \in S^2$ \\
Nonconformity score & $\dgeo(\hat y(x),y)/\hat\sigma(x)$ & $|y-\hat y(x)|/\hat\sigma(x)$ \\
Role of geodesic distance & Defines the score itself & Defines the neighbourhood for $\hat\sigma(x)$ \\
Prediction region & Geodesic cap on $S^2$, position-independent area & Real interval, location-dependent radius \\
\bottomrule
\end{tabular}
\end{table}

This method is not a direct re-implementation of \citet{amirishahbazi2026}'s algorithm. It is a structurally analogous method for a different problem class, which is closer in spirit to a spherical-input generalisation of ordinary locally-adaptive conformal prediction \citep{lei2018,romano2019}. The geodesic geometry is used on the input side to estimate local prediction difficulty. Therefore, it should not be described as a literal application of manifold-response adaptive geodesic conformal prediction. Besides, the vector/wind-field model from \citet{thomas2025}, discussed as future work in \cref{ch:conclusion}, provides a different setting.  In the vector-field case, the response itself has manifold structure, which makes the vector-field problem much closer to the setting considered by \citet{amirishahbazi2026}, and can potentially be applied more directly. This is because, after appropriate normalisation, the wind response can be represented as a direction on the sphere, allowing geodesic distance to be used between predicted and observed responses.
 
% ============================================================
